\section*{Chapter \ref{ch:b2:causality-and-experiments} --- Causality and the Experimental Ideal}

\begin{solution}{exo:b2:causality-and-experiments:1}
% figdata: means
Control 4,555 dollars (3,888 to 5,221); trained 6,349 (5,215 to 7,483); difference 1,794. The intervals overlap, by six
dollars, yet the difference is significant (Exercise 5): overlapping intervals for two means do not imply that their
difference is insignificant, since the interval of a difference is narrower than the sum of the two half-widths.
\end{solution}

\begin{solution}{exo:b2:causality-and-experiments:2}
The patients who go to hospital are sicker than those who stay home: their health without the hospital, $Y(0)$, is worse.
The \omterm{def:b2:causality-and-experiments:selection-bias}{selection bias} is negative, and a raw comparison of the health of patients and non-patients can make the hospital look
harmful even if it helps.
\end{solution}

\begin{solution}{exo:b2:causality-and-experiments:3}
Not by itself: with eight comparisons the chance of at least one significant difference is 0.34 under perfect
randomisation. Check how the randomisation was done, report the table, and show the estimate with the diploma (and the
other pre-treatment variables) as controls: 1,676 dollars, close to 1,794.
\end{solution}

\begin{solution}{exo:b2:causality-and-experiments:4}
Participants may find jobs partly by taking them from non-participants, an effect the local comparison does not see (a
failure of the \omterm{def:b2:causality-and-experiments:sutva}{stable unit treatment value assumption}); and the volunteers of one town may gain more than the population
at large.
\end{solution}

\begin{solution}{exo:b2:causality-and-experiments:5}
% code: listings/nsw_experiment.py
$1{,}794.3 \pm 1.96 \times 671.0$: from 479 to 3,110 dollars. Zero is outside it.
\end{solution}

\begin{solution}{exo:b2:causality-and-experiments:6}
% code: listings/selection_sim.py
ATE 1.001, ATT 1.284. The effect is $1 - 0.5 \times \mathit{skill}$ plus noise, and those who take the training have low
skill, so their mean effect is higher. If the more skilled took it, their untreated earnings would be higher than the
others': the \omterm{def:b2:causality-and-experiments:selection-bias}{selection bias} would be positive and the raw difference would overstate the effect, while the ATT would fall
below the ATE.
\end{solution}

\begin{solution}{exo:b2:causality-and-experiments:7}
% code: exercises/solutions.py low_earners
4,231 survey men earned less than 5,000 dollars in 1975. The raw difference becomes 127 dollars (s.e.\ 588) and the
adjusted one 1,193 (s.e.\ 654). The cut trims only the comparison group: it drops survey men with 1975 earnings above those of almost all
trainees, and keeps the 20 trainees above 5,000 dollars. That makes the comparison group more like the trained men, so the controls have less to do; the estimate moves toward the experimental 1,794, without
reaching it.
\end{solution}

\begin{solution}{exo:b2:causality-and-experiments:8}
In the experiment the controls are independent of treatment, so they change the estimate only through chance imbalances
such as the diploma. In the survey comparison the groups differ enormously in age, schooling, marriage and past earnings,
which also predict 1978 earnings: the \omterm{def:b2:multiple-regression:ovb}{omitted-variable bias} is large, and the controls remove most of it.
\end{solution}

\begin{solution}{exo:b2:causality-and-experiments:9}
Pupils who chose the programme may be more motivated; their parents may be more involved; their schools may have more
resources. An experiment: among pupils (or classes, to avoid spillovers) whose parents agree to take part, draw by lottery
who gets the programme, and compare scores after a year, with standard errors clustered by class if classes are drawn.
\end{solution}

\begin{solution}{exo:b2:causality-and-experiments:10}
Under random assignment of $n_1$ of $n$ units, each unit is treated with probability $n_1/n$, so $E[\bar Y_1] =
E[Y(1)]$ and $E[\bar Y_0] = E[Y(0)]$; the difference has expectation $E[Y(1) - Y(0)] = \mathit{ATE}$. With
independent draws from a large population, $\operatorname{Var}(\bar Y_1 - \bar Y_0) = \sigma_1^2/n_1 + \sigma_0^2/n_0$,
where $\sigma_d^2$ is the variance of $Y(d)$; it is estimated by the sample variances, as in the listing.
\end{solution}

\begin{solution}{exo:b2:causality-and-experiments:11}
Difference $= \mathit{ATE} + \text{selection bias} + (1 - p)(\mathit{ATT} - \mathit{ATU})$, with $p$ the share treated and
$\mathit{ATU}$ the effect on the untreated: since $\mathit{ATE} = p\,\mathit{ATT} + (1 - p)\mathit{ATU}$, $\mathit{ATT} =
\mathit{ATE} + (1 - p)(\mathit{ATT} - \mathit{ATU})$. The third term is the bias from heterogeneous effects: if those who
are treated gain more than those who are not, the difference overstates the ATE even without selection on $Y(0)$.
\end{solution}

\begin{solution}{exo:b2:causality-and-experiments:12}
Given $X_i = x$, treatment is independent of $(Y(0), Y(1))$, so $E[Y \mid D = 1, X = x] = E[Y(1) \mid X = x]$ and $E[Y
\mid D = 0, X = x] = E[Y(0) \mid X = x]$. Their difference is the effect for units with $X = x$. Averaging over the
distribution of $X$ gives $E[Y(1) - Y(0)] = \mathit{ATE}$ by the law of iterated expectations.
\end{solution}

\begin{solution}{pb:b2:causality-and-experiments:1}
% code: exercises/solutions.py population, rct, rollout, n_per_arm, power
\textbf{Part I.}
1. $Y_i(1) = 25{,}500 + 3{,}500m_i + 4{,}000\varepsilon_i$.
2. 1,500 dalens, since $E[m_i] = 0$.
3. Those with more drive.
4. 1,499 dalens.
5. Each worker is observed either with or without a voucher, never both.
6. Trained workers may take jobs that untrained ones would have had; vouchers given to many may raise the price of
courses.
\textbf{Part II.}
7. 1,100 dalens.
8. 362; from 389 to 1,810.
9. Yes: 1,499 is inside it.
10. Which workers got a voucher was drawn at random; another draw would give another estimate. The standard error measures
that spread.
11. Random assignment makes voucher holders a random half of the 800, so their mean effect equals the mean effect of all.
12. Leave the expected estimate unchanged and, since earlier earnings predict later earnings, shrink the standard error.
\textbf{Part III.}
13. 51 per cent.
14. 5,272 dalens.
15. 1,782 dalens.
16. 3,491 dalens: applicants have more drive, so they would have earned more anyway; $1{,}781.6 + 3{,}490.7 = 5{,}272.3$.
17. The effect rises with drive, and applicants have more of it.
18. Only if, among workers with the same education, applying were unrelated to drive: the \omterm{def:b2:causality-and-experiments:cia}{conditional independence
assumption}. Since drive is not measured, there is no reason to believe it.
\textbf{Part IV.}
19. $z_{0.025} = 1.96$, $z_{0.2} = 0.84$: $n = 2 \times 2.80^2\sigma^2/\delta^2 = 15.7\sigma^2/\delta^2$.
20. $15.7 \times 5{,}000^2/1{,}500^2 = 174.4$: 175 per arm.
21. With $\sigma = 5{,}217$: 190 per arm.
22. With $\sigma = 5{,}000$: $15.7 \times 5{,}000^2/750^2 = 698$ per arm.
23. $n$ is proportional to $1/\delta^2$.
24. Yes: with 400 per arm the power to detect 1,500 dalens is 0.98.
25. The pilot's estimate of the voucher's effect is 1,100 dalens a year (95 per cent interval 389 to 1,810), consistent with
the true 1,499; a roll-out to applicants would have shown 5,272, a bias of 3,491; 190 workers per arm suffice for 80 per
cent power against an effect of 1,500.
\end{solution}
